\begin{textsample}{POS Dim 1 – al – Score 38.00 – al/al\_ef\_21.txt}  \label{ex:f1_pos_160}
Reiteramos aqui, conforme já ressaltado na perspectiva da BNCC, que as habilidades não são desenvolvidas de forma descontextualizada, portanto, a proposta deste referencial engloba a leitura de textos que pertencem a gêneros que circulam nos mais diversos contextos de uso.
Apresentamos, a seguir, os campos de uso destacando as habilidades de leitura, oralidade e escrita, de forma contextualizada pelas práticas, gêneros e diferentes objetos do conhecimento em questão.
O \textbf{aprofundamento} cognitivo e a complexidade textual das \textbf{aulas} de leitura devem aumentar, progressivamente, desde os Anos Iniciais do Ensino Fundamental até o Ensino Médio.
Esta complexidade se expressa pela articulação : - da diversidade dos gêneros textuais escolhidos e das práticas consideradas em cada campo ; - da complexidade textual que se concretiza pela temática, estruturação sintática, vocabulário, recursos estilísticos utilizados, orquestração de \textbf{vozes} e linguagens presentes no texto ; - do uso de habilidades de leitura que \textbf{exigem} processos mentais necessários e progressivamente mais demandantes, passando de processos de recuperação de informação ( identificação, reconhecimento, organização ) a processos de \textbf{compreensão} ( comparação, distinção, estabelecimento de relações e inferência ) e de reflexão sobre o texto ( justificação, \textbf{análise}, articulação, apreciação e valorações estéticas, éticas, políticas e ideológicas ) ; - da consideração da cultura \textbf{digital} e das TDIC ; - da consideração da diversidade cultural, de \textbf{maneira} a abranger produções e formas de expressão diversas, a literatura infantil e juvenil, o cânone, o culto, o popular, a cultura de massa, a cultura das mídias, as culturas juvenis etc., de forma a garantir ampliação de repertório, além de interação e trato com o diferente.
É importante que os estudantes sejam partícipes em diversas atividades de leitura com níveis crescentes de complexidade, possibilitando que se amplie o seu repertório textual.
Diante disso, é possível que o professor possa trabalhar com gêneros que não estejam \textbf{sugeridos} em determinados anos.
Assim, se fizer mais sentido que um gênero mencionado e/ou habilidades a ele relacionadas no 9º ano sejam trabalhados no 8º, isso não \textbf{configura} um \textbf{problema}, desde que ao \textbf{final} do nível a diversidade indicada tenha sido contemplada.
Mesmo em relação à \textbf{progressão} das habilidades, seu desenvolvimento não se dá em curto espaço de tempo, podendo \textbf{supor} diferentes graus e ir se complexificando durante vários anos.
O Eixo da Produção de Textos busca contemplar as práticas de linguagem relacionadas à interação e à \textbf{autoria} ( individual ou coletiva ) do texto escrito, oral e \textbf{multissemiótico}, com diferentes finalidades e projetos enunciativos.
O trabalho com a produção de textos se fará a partir de gêneros a serem produzidos com enfoque nas mais variadas situações comunicativas comuns ao cotidiano dos estudantes.
Assim, contemplam -se produções, tais como as descritas abaixo : - construir um álbum de personagens famosas, de heróis/heroínas ou de vilões ou vilãs ; - produzir um almanaque que retrate as práticas culturais da comunidade ; narrar fatos cotidianos, de forma \textbf{crítica}, lírica ou bem-humorada em uma crônica ; - comentar e indicar diferentes produções culturais por meio de resenhas ou de playlists comentadas ; - descrever, avaliar e recomendar ( ou não ) um game em uma resenha, gameplay ou vlog ; escrever verbetes de curiosidades científicas ; - \textbf{sistematizar} dados de um estudo em um relatório ou relato multimidiático de campo ; - \textbf{divulgar} conhecimentos específicos por meio de um verbete de enciclopédia \textbf{digital} colaborativa ; - relatar fatos \textbf{relevantes} para a comunidade em notícias ; - cobrir acontecimentos ou levantar dados \textbf{relevantes} para a comunidade em uma reportagem ; - expressar posição em uma carta de leitor ou artigo de opinião ; - denunciar situações de desrespeito aos direitos por meio de fotorreportagem, fotodenúncia, poema, lambe-lambe, microrroteiro, dentre outros.
O tratamento das práticas de produção de textos compreende dimensões \textbf{inter-relacionadas} às práticas de uso e reflexão, tais como : - Consideração e reflexão sobre as condições de produção dos textos que regem a circulação de diferentes gêneros nas diferentes mídias e campos de atividade humana - Refletir sobre diferentes contextos e situações sociais em que se produzem textos e sobre as diferenças em termos formais, estilísticos e linguísticos que esses contextos determinam, incluindo -se aí a multissemiose e características da conectividade ( uso de hipertextos e hiperlinks, dentre outros, presentes nos textos que circulam em contexto \textbf{digital} ). - Analisar as condições de produção do texto no que \textbf{diz} respeito ao lugar social assumido e à imagem que se pretende passar a respeito de si mesmo ; ao leitor pretendido ; ao veículo ou à mídia em que o texto ou produção cultural vai circular ; ao contexto \textbf{imediato} e ao contexto \textbf{sócio-histórico} mais geral ; ao gênero do discurso/campo de atividade em questão etc. - Analisar aspectos sociodiscursivos, temáticos, composicionais e estilísticos dos gêneros propostos para a produção de textos, estabelecendo relações entre eles. - Dialogia e relação entre textos - Orquestrar as diferentes \textbf{vozes} nos textos pertencentes aos gêneros literários, fazendo uso adequado da “ fala ” do narrador, do discurso direto, indireto e indireto livre. - Estabelecer relações de intertextualidade para \textbf{explicitar}, sustentar e \textbf{qualificar} \textbf{posicionamentos}, construir e referendar explicações e relatos, fazendo usos de citações e paráfrases, devidamente marcadas e para produzir paródias e estilizações. - Alimentação temática - Selecionar informações e dados, \textbf{argumentos} e outras referências em fontes confiáveis impressas e \textbf{digitais}, organizando em roteiros ou outros formatos o material pesquisado, para que o texto a ser produzido tenha um nível de \textbf{aprofundamento} adequado ( para além do senso comum, quando for esse o caso ) e contemple a sustentação das posições defendidas. - Construção da textualidade - Estabelecer relações entre as partes do texto, levando em conta a construção composicional e o estilo do gênero, evitando repetições e usando adequadamente elementos coesivos que contribuam para a coerência, a continuidade do texto e sua \textbf{progressão} temática. - Organizar e/ou hierarquizar informações, tendo em \textbf{vista} as condições de produção e as relações lógico \textbf{discursivas} em jogo : causa/efeito ; tese/argumentos ; problema/solução ; definição/exemplos etc. - Usar recursos linguísticos e \textbf{multissemióticos} de forma articulada e adequada, tendo em \textbf{vista} o contexto de produção do texto, a construção composicional e o estilo do gênero e os efeitos de sentido pretendidos. - Aspectos notacionais e gramaticais - Utilizar, ao produzir textos, os conhecimentos dos aspectos notacionais – ortografia padrão, pontuação adequada, mecanismos de concordância nominal e verbal, regência verbal etc., sempre que o contexto \textbf{exigir} o uso da norma-padrão. - Estratégias de produção - Desenvolver estratégias de planejamento, revisão, \textbf{edição}, reescrita/redesign e avaliação de textos, considerando -se sua adequação aos contextos em que foram produzidos, ao modo ( escrito ou oral ; imagem estática ou em movimento etc. ), à variedade linguística e/ou semioses apropriadas a esse contexto, os enunciadores envolvidos, o gênero, o suporte, a esfera/campo de circulação, adequação à norma-padrão etc. - Utilizar softwares de \textbf{edição} de texto, de imagem e de áudio para editar textos produzidos em várias mídias, explorando os recursos multimídias disponíveis.
No eixo da produção deve -se conceber as práticas de produção buscando contemplar os gêneros que circulam nos mais diversos campos da atividade humana, priorizando aqueles que estejam presentes no cotidiano dos estudantes de nossa comunidade escolar.
Dessa forma, é preciso considerar, por exemplo, que nas atividades de escrita de um texto argumentativo no 7º ano, \textbf{devido} à \textbf{mobilização} frente a um tema, pode envolver \textbf{análise} e uso de diferentes tipos de \textbf{argumentos} e movimentos argumentativos, que podem estar previstos para o 9º ano.
Da mesma forma, o manuseio de uma ferramenta ou a produção de um tipo de \textbf{vídeo} proposto para uma apresentação oral no 9º ano pode se dar no 6º ou 7º anos, em função de um interesse que possa ter \textbf{mobilizado} os estudantes para \textbf{tanto}.

% matched lemmas: análise, aprofundamento, argumento, aula, autoria, compreensão, configurar, crítico, devido, digital, discursivo, divulgar, dizer, edição, exigir, explicitar, final, imediato, inter-relacionar, maneira, mobilizar, mobilização, multissemiótico, posicionamento, problema, progressão, qualificar, relevante, sistematizar, sugerir, supor, sócio-histórico, tanto, vista, voz, vídeo
\end{textsample}
